\chapter{基底与表面：器件真正开始的地方}

\begin{chaptergoal}
知道为什么 substrate 不只是机械支撑；理解介电常数、损耗、表面状态、平整度与清洗历史会进入谐振器和毫米波吸收问题。
\end{chaptergoal}

\section{选择基底时要问什么}

对微波谐振器，应关心介电常数、介质损耗、厚度及均匀性；对毫米/亚毫米波结构，还要考虑 substrate 中的传播、背短路或透镜几何。对加工，则要关心晶圆尺寸、平整度、表面粗糙度、是否有原生氧化层，以及后续沉积/刻蚀的兼容性。

\processchain{substrate / surface $\rightarrow$ 介电边界与界面状态 $\rightarrow$ 有效电容、场分布与损耗 $\rightarrow$ $f_0$、$Q_i$、optical absorption}

\section{“清洗”不是一个按钮}

清洗步骤应根据要去除的污染物和下一步工艺来定义。对于学习阶段，最重要的不是背化学配方，而是每次都明确：我要去掉什么？会不会留下新的表面状态？从清洗结束到沉积开始暴露了多久？这些问题应写进 process traveler。
